%% LyX 2.3.4.2 created this file.  For more info, see http://www.lyx.org/.
%% Do not edit unless you really know what you are doing.
\documentclass[11pt,spanish]{article}
\renewcommand{\rmdefault}{cmr}
\usepackage[T1]{fontenc}
\usepackage[utf8]{inputenc}
\usepackage{geometry}
\geometry{verbose,tmargin=2cm,bmargin=2cm,lmargin=2cm,rmargin=2cm}
\setlength{\parskip}{\bigskipamount}
\setlength{\parindent}{0pt}
\usepackage{amstext}
\usepackage{amsthm}
\usepackage{amssymb}

\makeatletter
%%%%%%%%%%%%%%%%%%%%%%%%%%%%%% Textclass specific LaTeX commands.
\theoremstyle{definition}
\newtheorem{defn}{\protect\definitionname}[section]
\theoremstyle{remark}
\newtheorem{rem}{\protect\remarkname}[section]
\theoremstyle{definition}
\newtheorem{xca}{\protect\exercisename}[section]
\theoremstyle{definition}
\newtheorem{example}{\protect\examplename}[section]

%%%%%%%%%%%%%%%%%%%%%%%%%%%%%% User specified LaTeX commands.
\usepackage{titling}
\setlength{\droptitle}{-2cm}
\usepackage{multicol}
\usepackage{enumitem}
\usepackage{proof}
\usepackage{mathrsfs}
\usepackage{amsfonts}
\usepackage{forest}
\usepackage{ebproof}
\usepackage{amssymb}
\definecolor{Guinda}{rgb}{0.49, 0.05, 0.12}
\definecolor{Rojo}{rgb}{0.8, 0, 0}
\definecolor{Verde}{rgb}{0,0.4,0}
\definecolor{Azul}{rgb}{0,0.4,0.4}
\definecolor{Naranja}{rgb}{0.8,0.4,0}
\usepackage{tikz}
\usetikzlibrary{automata, positioning, arrows}
\raggedcolumns

\makeatother

\usepackage{babel}
\addto\shorthandsspanish{\spanishdeactivate{~<>.}}

\usepackage{listings}
\lstset{mathescape=true}
\providecommand{\definitionname}{Definición}
\providecommand{\examplename}{Ejemplo}
\providecommand{\exercisename}{Ejercicio}
\providecommand{\remarkname}{Observación}

\begin{document}
\title{Teoría de la Computación, 2021-1\\
Nota de clase 6: Gramáticas Libres de Contexto}
\author{Manuel Soto Romero}
\date{\today \\ Colegio de Ciencia y Tecnología UACM}

\maketitle
\renewcommand*{\thefootnote}{\arabic{footnote}}
\setcounter{footnote}{0}
\setcounter{section}{6}
\renewcommand*\lstlistingname{C\'odigo}

\subsection{Introducción}

Como hemos visto, un mecanismo relevante para generar un lenguaje
es mediante el concepto de gramática formal. Las gramáticas formales
fueron introducidas por Noam Chomsky en 1956. La intención era tener
un modelo para la descripción de lenguajes naturales. Posteriormente
una clase especial de estos formalismos, llamados gramáticas libres
de contexto, se utilizaron como herramienta para presentar la sintaxis
de lenguajes de programación y para el diseño de analizadores léxicos
de compiladores.

\bigskip{}

\begin{defn}
(GLC) \emph{Una gramática libre de contexto (GLC) es una cuaterna
$G=\left\langle V,T,S,P\right\rangle $ tal que:}
\end{defn}
\begin{itemize}
\item \emph{V es un alfabeto de variables o símbolos no terminales, los
cuales se denotan por mayúsculas A, B, C, ...}
\item \emph{T es un alfabeto de constantes símbolos terminales, los cuales
se denotan por minúsculas a, b, c con la condición de que $T\cap V=\varnothing$.}
\item \emph{$S\in V$ es una variable distinguida llamada el símbolo inicial.}
\item $P$ \emph{es un conjunto finito de reglas de producción, donde cada
regla es de la forma}

\[
A\rightarrow\alpha
\]

y $A\in V,\alpha\in\left(V\cup T\right)^{*}$.
\end{itemize}
Veamos algunos ejemplos:
\begin{itemize}
\item Gramática que reconoce el lenguaje $a^{*}$:

\[
\begin{array}{rcl}
S & \rightarrow & aS\\
S & \rightarrow & \varepsilon
\end{array}
\]

\item Gramática que reconoce al lenguaje $a^{*}b^{*}$:

\[
\begin{array}{rcl}
S & \rightarrow & aS\\
S & \rightarrow & bA\\
S & \rightarrow & \varepsilon\\
A & \rightarrow & bA\\
A & \rightarrow & b\\
A & \rightarrow & \varepsilon
\end{array}
\]

\item Gramática que reconoce al lenguaje $0^{*}10^{*}10^{*}$:

\[
\begin{array}{rcl}
S & \rightarrow & A1A1A\\
A & \rightarrow & 0A\\
A & \rightarrow & \varepsilon
\end{array}
\]

\item Gramática que genere el lenguaje $L=\left\{ a^{i}b^{i}:i\geq0\right\} $:

\[
\begin{array}{rcl}
S & \rightarrow & aSb\\
S & \rightarrow & \varepsilon
\end{array}
\]

\end{itemize}
\bigskip{}

\begin{rem}
Es importante notar que toda gramática regular es libre de contexto
pero no al revés.
\end{rem}
%
\bigskip{}

\begin{xca}
Encontrar una GLC que genere el lenguaje de todos los palíndromos
sobre $\Sigma=\left\{ a,b\right\} $.
\end{xca}
\bigskip{}

\begin{xca}
Encontrar una GLC que genere el lenguaje de todas las cadenas sobre
$\Sigma=\left\{ a,b\right\} $ que tienen un número par de símbolos.
\end{xca}

\subsection{Gramáticas ambiguas}

Recordemos que dada una gramática podemos dar una derivación que permita
analizar la forma en que una cadena puede ser construida así como
analizar su respectivo árbol de derivación. En este caso, una derivación
usando gramáticas libres de contexto no suele ser única pues cada
paso depende de la variable elegida para su reescritura. Esto genera
un no determinismo que puede eliminarse usando una convención, la
cual suele ser de dos tipos, eligiendo siempre la variable más a la
izquierda o bien más a la derecha para el proceso de reescritura. 

\bigskip{}

\begin{defn}
\emph{(Derivación por la izquierda) Una derivación es por la izquierda
si en cada paso se reescribe la variable más a la izquierda en la
cadena.}
\end{defn}
\bigskip{}

\begin{defn}
\emph{(Derivación por la derecha) Una derivación es por la derecha
si en cada paso se reescribe la variable más a la derecha en la cadena.}
\end{defn}
%
\bigskip{}

\begin{example}
En la gramática
\end{example}
\[
\begin{array}{rcl}
S & \rightarrow & aSb\\
S & \rightarrow & a\\
A & \rightarrow & SbA\\
A & \rightarrow & SS\\
A & \rightarrow & ba
\end{array}
\]

\begin{enumerate}
\item La siguiente derivación por la izquierda de la cadena $aabbaa$:

\[
S\rightarrow aAS\rightarrow aSbAS\rightarrow aabAS\rightarrow aabbaS\rightarrow aabbaa
\]

\item La siguiente derivación por la derecha de la cadena $aabbaa$:

\[
S\rightarrow aAS\rightarrow aAa\rightarrow aSbAa\rightarrow aSbbaa\rightarrow aabbaa
\]

\end{enumerate}
%
\bigskip{}

\begin{defn}
\emph{(Gramática ambigua) Una GLC es ambigua si existe una cadena
para la cual existen dos derivaciones diferentes. De la misma forma,
podemos decir que es ambigua si existen dos árboles de derivación
diferentes para la misma cadena.}
\end{defn}
%
\bigskip{}

\begin{example}
Consideremos el alfabeto $\Sigma=\left\{ 0,1,+,*,(,)\right\} $. La
siguiente gramática para generar número naturales, sumas y productos,
en numeración binaria es ambigua, recordemos que el símbolo $|$ se
lee como ``o'' y permite simplificar las reglas de escritura:
\end{example}
\[
S\rightarrow S+S\;|\;S*S\;|\;\left(S\right)\;|\;0S\;|\;1S\;|\;0S\;|\;1S\;|\;0\;|\;1
\]

La cadena $1+1*0$ tiene dos derivaciones por la izquierda diferentes:

\[
S\rightarrow S+S\rightarrow1+S\rightarrow1+S*S\rightarrow1+1*S\rightarrow1+1*0
\]

\[
S\rightarrow S*S\rightarrow S+S*S\rightarrow1+S*S\rightarrow1+1*S\rightarrow1+1*0
\]

Los árboles de derivación correspondientes a la derivaciones anteriores
son:

\begin{multicols}{2}

\begin{center}
\begin{forest}	
	[S
		[S
			[1]
		]
		[+]
		[S
			[S
				[1]
			]
			[*]
			[S
				[0]
			]
		]
	]
\end{forest}
\end{center}

\columnbreak

\begin{center}
\begin{forest}	
	[S
		[S
			[S
				[1]
			]
			[+]
			[S
				[1]
			]
		]
		[*]
		[S
			[0]
		]
	]
\end{forest}
\end{center}

\end{multicols}

La ambigüedad se puede eliminar introduciendo paréntesis:

\[
S\rightarrow\left(S+S\right)\;|\;\left(S*S\right)\;|\;\left(S\right)\;|\;0S\;|\;1S\;|\;0S\;|\;1S\;|\;0\;|\;1
\]

Sin embargo, aunque la adición de los paréntesis elimina la ambigüedad,
las expresiones generadas tienen un excesivo número de paréntesis
lo que dificulta el análisis sintáctico (en un compilador por ejemplo).
Lo más común en estos casos es utilizar gramáticas como la original
siempre y cuando se establezca un orden de precedencia para los operadores.
Lo usual es establecer que {*} tenga una mayor precedencia que +.
Por ejemplo, la expresión 1{*}1+0 se interpreta como (1{*}1)+0 y la
expresión 10+11{*}110+1 se interpreta como 10+(11{*}110)+1.

\bigskip{}

\begin{xca}
Indica si la siguiente gramática es ambigua. Justifica tu respuesta.
\end{xca}
\[
\begin{array}{rcl}
S & \rightarrow & aSa\;|\;\varepsilon\\
A & \rightarrow & bA\;|\;\varepsilon
\end{array}
\]


\subsection{Normalización de gramáticas}

La normalización de gramáticas consiste en transformar todas las producciones
de una gramática de manera que tengan cierta forma sintáctica en particular.
Así la normalización de gramáticas libres de contexto es útil para
homogeneizar la forma de las producciones, además para optimizar los
procesos de derivación de cadenas. Con las \emph{Formas Normales},
se facilita una solución al problema de la pertenencia, es decir,
decidir si una cadena pertenece o no a un lenguaje:
\begin{center}
\emph{Dada una gramática G y una cadena w, ¿se cumple que $w\in L\left(G\right)$? }
\par\end{center}
\begin{itemize}
\item Si la palabra $w$ es generada por $G$, el proceso para la construcción
de un árbol de derivación terminará eventualmente.
\item En caso contrario \textbf{no} podemos saber cuando parar en la construcción
de un árbol de derivación.
\item Pero si la gramática está en cierta forma normal, sí es posible saber
cuando parar en la búsqueda de árboles de derivación para contestar
el problema de la pertenencia.
\end{itemize}
Para obtener gramáticas normales, son necesarias varias transformaciones
que simplificarán y reorganizarán las producciones.

\subsubsection{Eliminación de variables inútiles}

En una GLC puede haber dos tipos de variables inútiles: aquéllas que
nunca aparecen en el proceso de ninguna derivación y aquéllas que
no se pueden convertir en cadenas terminales.

\bigskip{}

\begin{defn}
\emph{(Variable alcanzable) Una variable A es alcanzable si existen
$u,v\in\left(V\cup\Sigma\right)^{*}$ tales que $S\Rightarrow^{*}uAv$.
La variable inicial S es alcanzable por definición.}
\end{defn}
%
\bigskip{}

\begin{defn}
\emph{(Variable terminal) Una variable A es terminal si existen $w\in\Sigma^{*}$
tales que $A\Rightarrow^{*}w$. En particular si $A\rightarrow\varepsilon$
es una producción, entonces A es terminal.}
\end{defn}
%
\textbf{Algoritmo para encontrar las variables terminales}

El siguiente algoritmo sirve para encontrar todas las variables terminales
de una GLC.

\textbf{TERM$_{1}:=\left\{ A\in V:\text{Existe una producción de la forma }A\rightarrow w,w\in\Sigma^{*}\right\} $.}\\
\textbf{TERM$_{i+1}:=$ TERM$_{i}\cup\left\{ A\in V:\exists\text{ producción }A\rightarrow w,w\in\left(\Sigma\cup TERM_{i}\right)^{*}\right\} $.}

De esta forma obtenemos una sucesión creciente de conjuntos de variables:

\[
TERM_{1}\subseteq TERM_{2}\subseteq TERM_{3}\subseteq...
\]

El algoritmo se presenta entonces como:
\begin{itemize}
\item Inicializar: \textbf{TERM $:=\left\{ A\in V:\text{Existe una producción de la forma }A\rightarrow w,w\in\Sigma^{*}\right\} $.}
\item Repetir: \textbf{TERM $:=$ TERM} \textbf{$\cup\left\{ A\in V:\exists\text{ producción }A\rightarrow w,w\in\left(\Sigma\cup TERM\right)^{*}\right\} $.}
\item Hasta: No se añaden nuevas variables a \textbf{TERM.}
\end{itemize}
\bigskip{}

\begin{example}
Encontrar las variables terminales de la siguiente gramática
\end{example}
\[
\begin{array}{rcl}
S & \rightarrow & ACD\;|\;bBd\;|\;ab\\
A & \rightarrow & aB\;|\;aA\;|C\\
B & \rightarrow & aDS\;|\;aB\\
C & \rightarrow & aCS\;|\;CB\;|\;CC\\
D & \rightarrow & bD\;|\;ba\\
E & \rightarrow & AB\;|\;aDb
\end{array}
\]

\emph{Solución}

\textbf{TERM$_{1}=\left\{ S,D\right\} $}\\
\textbf{TERM$_{2}=\left\{ S,D\right\} \cup\left\{ B,E\right\} =\left\{ S,D,B,E\right\} $}\\
\textbf{TERM$_{3}=\left\{ S,D,B,E\right\} \cup\left\{ A\right\} =\left\{ S,D,B,E,A\right\} $}\\
\textbf{TERM$_{4}=\left\{ S,D,B,E,A\right\} \cup\varnothing=\left\{ S,D,B,E,A\right\} $}

$\therefore$ El conjunto de variables no terminales es $\left\{ C\right\} $

\textbf{Algoritmo para encontrar las variables alcanzables}

El siguiente algoritmo sirve para encontrar todas las variables alcanzables
de una GLC.

\textbf{ALC$_{1}:=\left\{ S\right\} $.}\\
\textbf{ALC$_{i+1}:=$ ALC$_{i}\cup\left\{ A\in V:\exists\text{ producción }B\rightarrow uAv,B\in ALC_{i}\;u,v\in\left(V\cup\Sigma\right)^{*}\right\} $.}

De esta forma obtenemos una sucesión creciente de conjuntos de variables:

\[
ALC_{1}\subseteq ALC_{2}\subseteq ALC_{3}\subseteq...
\]

El algoritmo se presenta entonces como:
\begin{itemize}
\item Inicializar: \textbf{ALC $:=\left\{ S\right\} $.}
\item Repetir: \textbf{ALC} \textbf{$:=$ ALC $\cup\left\{ A\in V:\exists\text{ producción }B\rightarrow uAv,B\in ALC\;u,v\in\left(V\cup\Sigma\right)^{*}\right\} $..}
\item Hasta: No se añaden nuevas variables a \textbf{ALC.}
\end{itemize}
\bigskip{}

\begin{example}
Encontrar las variables terminales de la siguiente gramática
\end{example}
\[
\begin{array}{rcl}
S & \rightarrow & aS\;|\;AaB\;|\;ACS\\
A & \rightarrow & aS\;|\;AaB\;|\;AC\\
B & \rightarrow & bB\;|\;DB\;|\:BB\\
C & \rightarrow & aDa\;|\;ABD\;|\;ab\\
D & \rightarrow & aD\;|\;DD\;|\;ab\\
E & \rightarrow & FF\;|\;aa\\
F & \rightarrow & aE\;|\;EF
\end{array}
\]

\emph{Solución}

\textbf{ALC$_{1}=\left\{ S\right\} $}\\
\textbf{ALC$_{2}=\left\{ S\right\} \cup\left\{ A,B,C\right\} =\left\{ S,A,B,C\right\} $}\\
\textbf{ALC$_{3}=\left\{ S,A,B,C\right\} \cup\left\{ D\right\} =\left\{ S,A,B,C,D\right\} $}\\
\textbf{ALC$_{4}=\left\{ S,A,B,C,D\right\} \cup\varnothing=\left\{ S,A,B,C,D\right\} $}

$\therefore$ El conjunto de variables no alcanzables es $\left\{ E,F\right\} $

\subsubsection{Eliminación de producciones $\varepsilon$}

\bigskip{}

\begin{defn}
\emph{(Producción) Una producción de la forma $A\rightarrow\varepsilon$
se llama producción $\varepsilon$}
\end{defn}
%
\bigskip{}

\begin{defn}
\emph{(Variable anulable) Una variable A se llama anulable si $A\Rightarrow^{*}\varepsilon$.}
\end{defn}
%
\textbf{Algoritmo para encontrar las variables anulables}

El siguiente algoritmo sirve para encontrar todas las variables alcanzables
de una GLC.

\textbf{ANUL$_{1}:=\left\{ A\in V:A\rightarrow\varepsilon\text{ es una producción}\right\} $.}\\
\textbf{ANUL$_{i+1}:=$ ANUL$_{i}\cup\left\{ A\in V:\exists\text{ producción }A\rightarrow w,w\in\left(ANUL_{i}\right)^{*}\right\} $.}

De esta forma obtenemos una sucesión creciente de conjuntos de variables:

\[
ANUL_{1}\subseteq ANUL_{2}\subseteq ANUL_{3}\subseteq...
\]

El algoritmo se presenta entonces como:
\begin{itemize}
\item Inicializar: \textbf{ANUL $:=\left\{ A\in V:A\rightarrow\varepsilon\text{ es una producción}\right\} $.}
\item Repetir: \textbf{ANUL $:=$ ANUL} \textbf{$\cup\left\{ A\in V:\exists\text{ producción }A\rightarrow w,w\in\left(ANUL\right)^{*}\right\} $.}
\item Hasta: No se añaden nuevas variables a \textbf{ANUL.}
\end{itemize}
Una vez que haya sido determinado el conjunto \textbf{ANUL} de variables
anulables, por medio del algoritmo anterior, las producciones $\varepsilon$
se pueden eliminar (excepto $S\rightarrow\varepsilon$) añadiendo
nuevas producciones que simulen el efecto de las producciones $\varepsilon$
eliminadas. Más concretamente, por cada producción $A\rightarrow u$
se añaden producciones de la forma $A\rightarrow v$ obtenidas suprimiendo
de la cadena $u$ una, dos o más variable anulables presentes de todas
las formas posibles.

\bigskip{}

\begin{example}
Encontrar las producciones $\varepsilon$ de la siguiente gramática
\end{example}
\[
\begin{array}{rcl}
S & \rightarrow & AB\;|\;ACA\;|\;ab\\
A & \rightarrow & aAa\;|\;B\;|\;CD\\
B & \rightarrow & bB\;|\;bA\\
C & \rightarrow & cC\;|\;\varepsilon\\
D & \rightarrow & aDc\;|\;CC\;|\;ABb
\end{array}
\]

\emph{Solución}

\textbf{ANUL$_{1}=\left\{ C\right\} $}\\
\textbf{ANUL$_{2}=\left\{ C\right\} \cup\left\{ D\right\} =\left\{ C,D\right\} $}\\
\textbf{ANUL$_{3}=\left\{ C,D\right\} \cup\left\{ A\right\} =\left\{ C,D,A\right\} $}\\
\textbf{ANUL$_{4}=\left\{ C,D,A\right\} \cup\left\{ S\right\} =\left\{ C,D,A,S\right\} $}\\
\textbf{ANUL$_{5}=\left\{ C,D,A,S\right\} \cup\varnothing=\left\{ C,D,A,S\right\} $}

$\therefore$ El conjunto de variables anulables es $\left\{ C,D,A,S\right\} $.
Al eliminar las producciones $\varepsilon$, tenemos:

\[
\begin{array}{rcl}
S & \rightarrow & AB\;|\;ACA\;|\;ab\;|\;B\;|\;CA\;|\;AA\;|\;AC\;|\;A\;|\;C\;|\;\varepsilon\\
A & \rightarrow & aAa\;|\;B\;|\;CD\;|\;aa\;|\;C\;|\;D\\
B & \rightarrow & bB\;|\;bA\;|\;b\\
C & \rightarrow & cC\;|\;c\\
D & \rightarrow & aDc\;|\;CC\;|\;ABb\;|\;ac\;|\;C\;|\;Bb
\end{array}
\]


\subsubsection{Eliminación de producciones unitarias}

\bigskip{}

\begin{defn}
\emph{(Producción unitaria) Una producción de la forma $A\rightarrow B$
donde $A$ y $B$ son variables, se llama producción unitaria.}
\end{defn}
%
\bigskip{}

\begin{defn}
\emph{(Conjunto unitario) El conjunto unitario de una variable A se
define de la siguiente manera:}

\[
UNIT(A):=\left\{ X\in V:\exists\text{ una derivación }A\Rightarrow^{*}X\text{ que usa únicamente producciones unitarias}\right\} 
\]

Por definición $A\in UNIT\left(A\right)$
\end{defn}
%
\textbf{Algoritmo para encontrar las producciones unitarias}

El siguiente algoritmo sirve para encontrar todas las variables alcanzables
de una GLC.

\textbf{UNIT$_{1}(A):=\left\{ A\right\} $.}\\
\textbf{UNIT$_{i+1}:=$ UNIT$_{i}(A)\cup\left\{ X\in V:\exists\text{ producción }Y\rightarrow X,Y\in\left(UNIT_{i}(A)\right)\right\} $.}

De esta forma obtenemos una sucesión creciente de conjuntos de variables:

\[
UNIT_{1}\subseteq UNIT_{2}\subseteq UNIT_{3}\subseteq...
\]

El algoritmo se presenta entonces como:
\begin{itemize}
\item Inicializar: \textbf{UNIT$(A):=\left\{ A\right\} $.}
\item Repetir: \textbf{UNIT$(A):=$ UNIT$(A)\cup\left\{ X\in V:\exists\text{ producción }Y\rightarrow X,Y\in\left(UNIT(A)\right)\right\} $}
\item Hasta: No se añaden nuevas variables a \textbf{ANUL.}
\end{itemize}
Las producciones unitarias se pueden eliminar añadiendo para cada
variable $A$ las producciones (no unitarias) de las variables contenidas
en el conjunto unitario UNIT(A). 

\bigskip{}

\begin{example}
Eliminar las producciones unitarias de la siguiente gramática:
\end{example}
\[
\begin{array}{rcl}
S & \rightarrow & AS\;|\;AA\;|\;BA\;|\;\varepsilon\\
A & \rightarrow & aA\;|\;a\\
B & \rightarrow & bB\;|\;bC\;|\;C\\
C & \rightarrow & aA\;|\;bA\;|\;B\;|ab
\end{array}
\]

\emph{Solución}

\textbf{UNIT(S)$_{1}=\left\{ S\right\} $}\\
\textbf{UNIT(S)$_{2}=\left\{ S\right\} \cup\varnothing=\left\{ S\right\} $}\\
\textbf{UNIT(A)$_{1}=\left\{ A\right\} $}\\
\textbf{UNIT(A)$_{2}=\left\{ A\right\} \cup\varnothing=\left\{ A\right\} $}\\
\textbf{UNIT(B)$_{1}=\left\{ B\right\} $}\\
\textbf{UNIT(B)$_{2}=\left\{ B\right\} \cup\left\{ C\right\} =\left\{ B,C\right\} $}\\
\textbf{UNIT(B)$_{3}=\left\{ B,C\right\} \cup\left\{ B\right\} =\left\{ B,C\right\} $}\\
\textbf{UNIT(C)$_{1}=\left\{ C\right\} $}\\
\textbf{UNIT(C)$_{2}=\left\{ C\right\} \cup\left\{ B\right\} =\left\{ C,B\right\} $}\\
\textbf{UNIT(C)$_{3}=\left\{ C,B\right\} \cup\left\{ C\right\} =\left\{ C,B\right\} $}

Eliminando las producciones unitarias se obtiene:

\[
\begin{array}{rcl}
S & \rightarrow & AS\;|\;AA\;|\;BA\;|\;\varepsilon\\
A & \rightarrow & aA\;|\;a\\
B & \rightarrow & bB\;|\;bC\;|\;aA\;|\;bA\;|\;ab\\
C & \rightarrow & aA\;|\;bA\;|\;bB\;|\;bC\;|ab
\end{array}
\]

\begin{thebibliography}{1}
\bibitem{key-1}Favio E. Miranda, A. Liliana Reyes, et.al, \emph{Notas
de Clase para el curso de Autómatas y Lenguajes Formales}, Facultad
de Ciencias, UNAM, Revisión 2019-2.

\bibitem{key-2}Rodrigo De Castro, \emph{Notas de Clase de Teoría
de la Computación: Lenguajes, autómatas, gramáticas,} Universidad
Nacional de Colombia, Facultas de Ciencias, 2004.
\end{thebibliography}

\end{document}
